\section{Methods}
\label{sec:methods}

\subsection{CENSO: binding the limit to the result}
\label{sec:methods-vocab}
CENSO\footnote{\url{https://w3id.org/censo/}, CC~BY~4.0.} extends
SOSA/SSN~\cite{janowicz2019sosa,haller2019ssn} for observations from a method
with a detection limit, reusing PROV-O~\cite{w3c2013provo} for the analytical
run, SKOS~\cite{w3c2009skos} for substance concepts and QUDT for units;
substances are aligned to ChEBI~\cite{degtyarenko2008chebi} where an identifier
exists (\num{180} of \num{240}). It follows LOT~\cite{poveda2022lot} (successor
to NeOn~\cite{suarez2012neon}), whose competency questions produced the
evaluation in Section~\ref{sec:results-cq}, and is published for
reuse~\cite{wilkinson2016fair}. Four commitments distinguish it from the
vocabularies surveyed in Section~\ref{sec:results-gap}.

\paragraph{The limit is carried onto the result.}
\texttt{censo:AnalyticalMethod} is a \texttt{sosa:Procedure} carrying a unit and
at least one limit --- of quantification, of detection, or both, since a
constraint satisfiable only by inventing the missing datum constrains nothing.
That much CHMO also has, and puts in the same place
(Section~\ref{sec:results-gap}). The step that separates a bound from a
measurement is the next one: \texttt{censo:resultUpperBound} places the limit
on the \emph{observation}, so \emph{censored} is a property of a result and
not of the procedure that produced it. Two observations from one method, one
below the limit and one above, differ in exactly this.

An earlier version of the vocabulary reified the analytical run for the same
purpose. It is retired, for two reasons that are worth stating because the
commitment is unchanged. A run is a \texttt{prov:Activity} that used a sample
and generated an observation, and the method it executed already carries the
limits, so declaring our own terms for it reinvented \textsc{prov-o}. And no
European release reports a run identifier: an annual mean spans many runs, so
minting one per station-year would not simplify but misstate. Where a publisher
does have run-level provenance, \textsc{prov} expresses it with no extension,
and the bound it produces still lands on
\texttt{censo:resultUpperBound}.\src{scripts/99\_audit.py, \texttt{check\_no\_dead\_terms}}

\paragraph{The substances are reconciled, not renamed.}
A \texttt{censo:Analyte} is a \texttt{sosa:ObservableProperty} --- the
concentration \emph{of} a substance --- and not the substance, whose instances
are molecules. The two are related by \texttt{rdfs:seeAlso} to the substance's
ChEBI class, in a separate module the vocabulary does not
import\src{scripts/20\_align\_external.py}: \num{299} of \num{354} analytes
across the two packages resolve to exactly one ChEBI class by CAS registry
number, \num{26} resolve to more than one and are deliberately left unaligned,
and \num{29} --- technical mixtures and congener groups such as C10--13
chloroalkanes --- have no ChEBI entry at all, which is a fact about what the
regulation names rather than a gap in the mapping.

\texttt{owl:sameAs} is reserved for the relation that is an identity: two
packages' names for one observable property. \num{30} such pairs reconcile the
European and Turkish vocabularies, so the co-regulated stratum of
Section~\ref{sec:results-dual} is an entailment rather than a string join, and
the audit requires the two to agree --- every substance the analysis treats as
co-regulated must be reachable through a \texttt{sameAs} pair.

\paragraph{Detection status is explicit and exhaustive.}
\texttt{CensoredObservation}, \texttt{EstimatedObservation},
\texttt{QuantifiedObservation} and \texttt{UnresolvedObservation} are pairwise
disjoint and cover \texttt{AssessedObservation}. Never measured is a fifth
situation, carried by the absence of an observation rather than by a class:
under the open-world assumption that is how one says a measurement was not made.
It is the distinction a \texttt{WHERE concentration = 0} query cannot draw, and
conflating it manufactures findings.

\paragraph{A threshold is a conditional judgement.}
\texttt{censo:Threshold} carries a value, a unit, a citation to the regulation,
and \texttt{requiresCondition} links to \texttt{ApplicabilityCondition}
subclasses: matrix, analytical fraction, bioavailability basis, hardness class.
Where a precondition is unmet the outcome is \texttt{PreconditionUnmet}, one
of the four subclasses of \texttt{IndeterminateCompliance} --- not a default
pass. Annex~I
creates two such conditions in footnotes: the standards for lead and nickel
``refer to bioavailable concentrations of the substances'', and cadmium's
``vary depending on the hardness of the water as specified in five class
categories''. The released package attaches them, and
Section~\ref{sec:results-precondition} shows what they decide.

\paragraph{Compliance is three-valued, and the third value records why.}
\texttt{Exceedance}, \texttt{Compliant} and \texttt{IndeterminateCompliance} are
jointly exhaustive \emph{for a given observation--threshold pair}, and they are
what the Directive itself admits: the standard is met, it is not met, or
compliance cannot be determined. Why it cannot is carried by four subclasses of
the third --- \texttt{PreconditionUnmet}, \texttt{MethodInsufficient},
\texttt{PossibleExceedance} and \texttt{BoundNotEstablished} --- and not by a
fourth top-level value, which would be the vocabulary asserting a distinction
the Directive does not make. \texttt{PossibleExceedance} belongs there rather
than beside \texttt{Exceedance}: an interval straddling the limit decides
nothing, which is indeterminacy with a particular cause, and the assessment
reported in Section~\ref{sec:results-wbabox} has counted it that way throughout.

Where two subclasses describe one result --- a censored result whose limit
exceeds the standard has an interval straddling it too --- Article~3(3b) is
decisive and the observation is recorded under \texttt{MethodInsufficient}. That
precedence is a rule-layer priority and deliberately not an
\texttt{owl:disjointWith} axiom between the subclasses, for the same reason no
class-level disjointness is asserted among the three outcomes: a measurement may
legitimately be compliant under one jurisdiction's standard and indeterminate
under another's, and Section~\ref{sec:results-dual} depends on the graph holding
both. Exclusivity is carried instead by an \texttt{owl:AllDisjointProperties}
axiom over \texttt{exceeds}, \texttt{possiblyExceeds} and
\texttt{belowThreshold}, with the equivalent pairwise axioms alongside it
because the pure-Python engine used here implements only the pairwise rule.

Figure~\ref{fig:geometry} shows three real observations. Where the threshold
clears the censored interval (a), the non-detection decides the question and the
ontology returns \texttt{Compliant}: censoring is not an excuse for refusing
every verdict. Where it falls inside (b), no result the method could have
returned would have decided the question, and any substitution rule merely picks
one.

The panels are cut at the quantification limit: Waterbase reports
\texttt{procedureLOQValue} and no detection limit, and both provisions at issue
--- Article~3(3b) of Directive 2008/105/EC and Article~4(1) of Directive
2009/90/EC --- are written about that limit. Section~\ref{sec:methods-noise}
distinguishes a threshold below the detection limit from one between the two;
this source cannot show that on measured cases.

\begin{figure}[htbp]\centering
\includegraphics[width=\textwidth]{fig06_decision_geometry.pdf}
\caption{Why two verdicts are not enough, on three real observations. A
non-detection is the interval $[0,\mathrm{LOQ}]$ (shaded), not a value; the
threshold $T$ is a line. Panels (a) and (c) are decidable and a conventional
pipeline gets them right; panel (b) is not, and it is silently recorded as
compliant --- a distinction available to a reasoner but invisible to a numeric
column. Each panel is scaled to its own limit; substance, country and values are
in \texttt{supplementary/figure\_data/fig06\_decision\_geometry.csv}. Generated by
\texttt{scripts/90\_figures.py}.}
\label{fig:geometry}
\end{figure}

\subsection{Regulations as pluggable, versioned packages}
\label{sec:methods-packages}
Thresholds expire. \texttt{cereg:RegulationPackage} makes each regulation a
self-contained module of individuals with jurisdiction, dates and a
machine-readable \texttt{transcriptionStatus}, so unverified values can be
excluded rather than trusted. Two packages are released with the
ontology\src{scripts/19\_build\_regulation\_packages.py}: \texttt{tr-ysky-2016}
(\num{434} thresholds) and \texttt{eu-2008-105-2026} (\num{114} thresholds),
both transcribed from the primary legal texts and carrying
\texttt{VerifiedAgainstPrimarySource}. Each contributes individuals only --- no
class, no property --- so one loads safely alongside another; the validator
checks this. Their thresholds carry different IRIs, so the Turkish and European
assessments coexist in one graph instead of overwriting each other.

\subsection{Is this not simply what detection limits are for?}
\label{sec:methods-noise}
Limits of detection exist to separate signal from noise, and a laboratory
reporting ``not detected'' is doing the right thing; nothing here disputes the
analytical practice. Two statements are then conflated:

\begin{itemize}
  \item \emph{the measurement cannot resolve a value} -- true below the limit;
    and
  \item \emph{the concentration is zero} -- not established by the measurement.
\end{itemize}

The laboratory asserts the first, the database records the second. What is lost
is not a value but the \emph{bound}: a non-detection establishes $c \in
[0,\mathrm{LOD}]$, and \num{0.0} is a different claim, not a weaker one. The
method is chosen on analytical grounds, the threshold arrives later from a legal
instrument, and their relative position is accidental. Article~3(3b) of the
consolidated Directive 2008/105/EC provides that a result reported as less than
a limit of quantification that is itself above the EQS ``shall not be considered
for the purposes of assessing the overall chemical status of that water
body''~\cite{eu2008eqs,eu2009qaqc}. The law requires a third outcome and
requires it to be identifiable; a two-valued schema cannot represent it.
Section~\ref{sec:results-waterbase} reports how often it arises.

\subsection{Could a column have done this?}
\label{sec:methods-why}
A relational table with a \texttt{censored} flag and an \texttt{lod} column
yields every number in Section~\ref{sec:results-waterbase}; our analysis scripts
are plain Python over CSV, and we do not claim these quantities are otherwise
uncomputable. What the flag does not give:

\begin{enumerate}
  \item \textbf{Several regulations over one set of observations.} The Turkish
    and European assessments differ substantially
    (Section~\ref{sec:results-waterbase}); with packages that is a re-query,
    with a threshold column a re-implementation, and the two cannot be held at
    once.
  \item \textbf{Preconditions belong to the threshold, not the measurement.} A
    flag on the observation cannot express a condition on the standard.
  \item \textbf{Justification.} A defined class yields a derivation for each
    verdict; a boolean yields a value.
  \item \textbf{Transfer.} The pattern is domain-independent; a column in one
    team's spreadsheet is not.
\end{enumerate}

This paper demonstrates (1) and (2) on real data and (3) on the axiom suite;
(4) is design intent, evidenced only by the vocabulary's independence from the
water domain, and we mark it as such.

\subsection{What OWL does, and what it does not}
OWL supplies identity, taxonomy, disjointness, classification and explanation.
It cannot compare data values, so \texttt{exceeds}, \texttt{possiblyExceeds} and
\texttt{belowThreshold} are materialised by the rule layer, then classified. The
open-world assumption that makes CENSO right for censoring blocks inferring
``none of the above'', so \texttt{IndeterminateCompliance} is materialised under
a locally closed world in SHACL. We do not claim that a description-logic
reasoner computes compliance.

\subsection{What is counted, and what is inferred}
\label{sec:methods-inference}
Almost nothing here is an estimate. The shares in
Section~\ref{sec:results-waterbase} and the three-valued assessment in
Section~\ref{sec:results-wbabox} are counts over every river station-year in the
release, computed by arithmetic on four reported fields; there is no sampling
frame and no population parameter to estimate, so no interval is reported for
them. Two places depart from that and are treated accordingly. Where a quantity
is selected \emph{by} its own value --- the ranked substances of
Figure~\ref{fig:loqeqs}a --- the \num{95}\,\% Wilson interval is plotted,
because ranking puts the least certain estimates at the top. Where a proportion
is compared between groups, Section~\ref{sec:results-dual} reports the interval
for the comparison it rests on.

Two limits of that treatment are worth stating rather than leaving to a reader.
Station-years are clustered --- by station, substance, reporting country and year
--- so they are not independent observations, and no comparison here corrects for
that; the descriptive shares do not need it, and the group comparisons should be
read as differences between the populations counted rather than as tests. And the
denominator of the two legal criteria is itself a reported quantity: a
station-year with no limit cannot fail a test about limits, so every rate built
on them is a lower bound, and loosest where reporting is worst
(Figure~\ref{fig:country-practice}).

\subsection{Reproducibility}
\label{sec:methods-repro}
Every number in this paper is produced by a script and written under
\texttt{eval/}; none is typed by hand. The pipeline is deterministic and runs in
pure Python (\texttt{owlrl}, \texttt{pyshacl}), so no external reasoner is
required. The ontology was scanned with OOPS!~\cite{poveda2014oops};
FOOPS!\ \cite{garijo2021foops} is the FAIR assessment we report as not yet run
rather than assert. A generated manifest maps every numeric claim to its
artefact and script, and fails if one cannot be
traced\src{scripts/95\_numbers\_manifest.py}.

The vocabulary is scanned for modelling pitfalls with
OOPS!~\cite{poveda2014oops} and assessed for FAIR publication with
FOOPS!~\cite{garijo2021foops}. FOOPS!\ is run against the published IRI rather
than the local file, because content negotiation, version-IRI resolution and
registry presence are properties of the publication and not of a serialisation:
it returns \num{0.78}, passing \num{18} of \num{24}
checks\src{scripts/15\_foops.py}. The six that do not pass are of two kinds and
we separate them rather than report a score. Three are missing annotations --- a
citation, a publisher, an issue date --- properties of the file, repairable
without touching an axiom. Three are absence from a public registry, which no
edit to the file achieves. Neither kind bears on the modelling claims made here;
the full report is in \texttt{eval/foops\_assessment.md} and the cached service
response beside it, so the figure can be re-derived without a network.
